\documentclass[useAMS,usenatbib]{mn2e}

% ─── PAQUETES ─────────────────────────────────────────────────────────────────
\usepackage[utf8]{inputenc}
\usepackage[T1]{fontenc}
\usepackage{amsmath,amssymb}
\usepackage{graphicx}
\usepackage{booktabs}
\usepackage{array}
\usepackage{multirow}
\usepackage{xcolor}
\usepackage{hyperref}
\usepackage{url}
\usepackage{siunitx}
\usepackage{textcomp}

% \doi no definido en mn2e — definición con hipervínculo
\newcommand{\doi}[1]{\href{https://doi.org/#1}{\nolinkurl{doi:#1}}}

% ─── METADATOS DE REVISTA ─────────────────────────────────────────────────────
\title[Bosque Oscuro Monte Carlo]{
  Del Silencio a la Dispersión:\\
  El Efecto Espacial Neto del Bosque Oscuro\\
  es Modesto y Crece con la Densidad
}

\author[Juan Galaz]{
  Juan Galaz$^{1}$\thanks{Correo electrónico: juan.galaz@proton.me;
  ORCID: \href{https://orcid.org/0009-0007-7474-7560}{0009-0007-7474-7560}}
  \\
  $^{1}$Investigador independiente, autofinanciado
}

\pubyear{2026}

\begin{document}

\maketitle

% ─── RESUMEN ──────────────────────────────────────────────────────────────────
\begin{abstract}
Presentamos una simulación Monte Carlo estocástica de la hipótesis del
Bosque Oscuro, la conjetura de que las civilizaciones galácticas se
destruyen mutuamente de forma preventiva al detectarse, y
caracterizamos las tasas de destrucción resultantes y las distribuciones
espaciales en dos regímenes de población.
Esta versión corrige un defecto del modelo previo (v5): la tasa de
falsos positivos $p_\text{fp}$ actuaba como piso de la probabilidad de
detección a cualquier distancia, generando más del 96 por ciento de las
destrucciones a distancias no físicas (mediana $\sim 32{,}000$\,al,
$32\times$ el radio máximo de detección justificable), y el 30 por
ciento de las destrucciones violaba causalidad.
El modelo corregido excluye $p_\text{fp}$ del canal de detección de
blancos reales e impone la condición causal
$d \leq c\,(t - t_\text{nac})$.
Bajo parámetros de Drake \emph{pesimistas} ($N_c = 440$, 50 semillas
independientes) el modelo corregido produce una tasa de destrucción de
$0{,}59 \pm 0{,}34$ por ciento, con supervivientes a una distancia media
al vecino más cercano de $1733 \pm 53$\,al y cociente
$r_\text{obs/P} = 1{,}42 \pm 0{,}04$ (superpoissoniano), consistente
con la concentración geométrica de nacimientos en la ZGH.
Simulaciones de control apareadas con $P_\text{ataque} = 0$ (50
semillas) producen $r_\text{obs/P}^\text{ctrl} = 1{,}41 \pm 0{,}04$;
el efecto diferencial neto del Bosque Oscuro es
$\Delta r = +0{,}0045 \pm 0{,}0032$: positivo pero minúsculo, la
estructura espacial en baja densidad está dominada por el perfil ZGH.
Bajo parámetros \emph{optimistas} ($N_c = 10{,}000$ acotado, 10
semillas) la tasa de destrucción asciende a $22{,}9 \pm 0{,}5$ por
ciento, con destrucciones a una mediana de $2{,}400$\,al y
supervivientes a $553 \pm 4$\,al del vecino más cercano.
El cociente obs/Poisson es $1{,}177 \pm 0{,}008$, frente al control
$r_\text{obs/P}^\text{ctrl} = 1{,}117 \pm 0{,}001$ ($n = 3$ semillas):
el efecto diferencial es $\Delta r_\text{opt} = +0{,}060 \pm 0{,}008$
($7{,}8\sigma$).
El Bosque Oscuro no concentra a los supervivientes: los \emph{dispersa},
de forma modesta y monotónica con la densidad
($+0{,}005 \to +0{,}060$), porque las destrucciones son locales
y eliminan preferentemente al vecino más cercano.
El resultado contrasta con la concentración masiva
($\Delta r = -0{,}54$) reportada por la versión previa del modelo, que
demostramos era atribuible en su totalidad al piso de falsos positivos:
un análisis de sensibilidad de un factor muestra que $p_\text{fp}$
era el parámetro individual dominante del modelo v5, no incluido en su
análisis de Sobol.
Estos hallazgos acotan el efecto espacial neto de la dinámica del
Bosque Oscuro a una señal de dispersión débil, detectable solo en
regímenes de alta densidad poblacional, y subrayan la necesidad de
verificar la coherencia física de los términos de ruido de detección
en modelos de la paradoja de Fermi.
\end{abstract}

\begin{keywords}
astrobiología --- inteligencia extraterrestre --- paradoja de Fermi ---
métodos Monte Carlo --- estadística espacial --- zona galáctica habitable
\end{keywords}

% ─────────────────────────────────────────────────────────────────────────────
\section{Introducción}
\label{sec:intro}
% ─────────────────────────────────────────────────────────────────────────────

La paradoja de Fermi, la aparente contradicción entre la alta
probabilidad a priori de civilizaciones tecnológicas extraterrestres y
su ausencia observacional, ha motivado una amplia clase de hipótesis
explicativas \citep{fermi1950,hart1975,brin1983,webb2002}.
Entre las más radicales se encuentra la \emph{hipótesis del Bosque
Oscuro} (HBO), en la que el silencio del cosmos es un equilibrio
estratégico: cualquier civilización que revela su ubicación arriesga
la aniquilación por un vecino más avanzado, por lo que la estrategia
óptima es el silencio o el ataque preventivo
\citep{cixin2008,devladar2013}.

Los modelos formales de la HBO siguen siendo escasos en la literatura
astrobiológica revisada por pares.
\citet{forgan2009} desarrolló el banco de pruebas numérico más citado
para hipótesis de la paradoja de Fermi, utilizando realizaciones Monte
Carlo de la ecuación de Drake, pero modeló la destrucción como una
probabilidad fija ($P_\text{destroy} = 0{,}5$) independiente de la
detección, la distancia o la disparidad tecnológica.
\citet{devladar2013} formalizó el fundamento de teoría de juegos de la
HBO, demostrando que el ataque preventivo es un equilibrio de Nash bajo
ciertas matrices de pagos, pero no simuló las consecuencias espaciales
a escala galáctica.
\citet{foucher2023} y \citet{forgan2011} señalaron independientemente
que la distribución espacial de los supervivientes debería depender de
la densidad de población, pero ninguno cuantificó la transición de
régimen.

El presente trabajo aborda esta brecha implementando un pipeline
Monte Carlo completo de detección--decisión--destrucción con los
siguientes elementos novedosos:
\begin{enumerate}
  \item Una función de probabilidad de detección tridimensional
        continua basada en la edad tecnológica y un radio máximo de
        detección físicamente motivado
        \citep[$r_\text{max} = 1000$\,al,][]{troitskij1989,loeb2007},
        con una condición explícita de causalidad
        ($d \leq c\,(t - t_\text{nac})$): ninguna civilización puede
        ser detectada antes de que su primera señal alcance al
        observador;
  \item Tipos sociológicos heterogéneos (agresivo, defensivo, pacifista,
        curioso) con umbrales de ataque dependientes del tipo, siguiendo
        a \citet{devladar2013};
  \item Un perfil de densidad continuo de la Zona Galáctica Habitable
        (ZGH) basado en \citet{lineweaver2004}, en reemplazo de la
        aproximación de anillo uniforme;
  \item Una memoria histórica de civilizaciones destruidas que reduce la
        probabilidad de comunicación en regiones donde se han observado
        destrucciones previas;
  \item Análisis de doble escenario (parámetros de Drake
        pesimistas/optimistas) que permite la comparación de regímenes;
  \item Un análisis de sensibilidad centrado en la coherencia física
        del canal de detección, motivado por la corrección de un
        artefacto dominante en la versión previa del modelo
        (\S\ref{sec:sensibilidad}, Apéndice~\ref{app:correction}).
\end{enumerate}

El resultado central es que el efecto espacial neto de la dinámica del
Bosque Oscuro es una \emph{dispersión} modesta de los supervivientes,
que crece monotónicamente con la densidad poblacional.
Para aislar este efecto de la concentración geométrica de la ZGH, el
trabajo incluye simulaciones de control con dinámica de destrucción
desactivada (Sección~\ref{sec:control}), cuyo contraste apareado
permite cuantificar el efecto diferencial neto del Bosque Oscuro
($\Delta r$).
Una versión previa de este modelo (v5) reportaba el resultado opuesto
—concentración masiva en el régimen denso, $\Delta r = -0{,}54$—;
en la Sección~\ref{sec:sensibilidad} y el
Apéndice~\ref{app:correction} demostramos que aquella señal era un
artefacto del término de falsos positivos, que actuaba como piso de
detección a cualquier distancia y dominaba la dinámica de destrucción.

Este modelo debe entenderse como una exploración del espacio de
parámetros de la HBO, no como una predicción cuantitativa del estado
actual de la galaxia.
Cuatro parámetros carecen de restricción observacional directa
($\theta_\text{threat}$, $d_0$, las fracciones sociológicas,
$p_\text{fn}$); los valores absolutos de las tasas de destrucción
dependen de ellos.
La lección metodológica de la corrección v5\,$\to$\,v6 es directa: la
coherencia física de los términos de ruido de detección debe
verificarse explícitamente —por ejemplo, comprobando que la
distribución de distancias de destrucción sea consistente con la
escala física del modelo— antes de interpretar señales espaciales.
El código fuente y los datos de salida están depositados en un
repositorio público para facilitar la reproducción y extensión del
modelo.

% ─────────────────────────────────────────────────────────────────────────────
\section{Modelo}
\label{sec:model}
% ─────────────────────────────────────────────────────────────────────────────

\subsection{Generación de civilizaciones}
\label{sec:civgen}

Las civilizaciones se generan muestreando la ecuación de Drake:
\begin{equation}
  N_c = N_\star \cdot f_p \cdot n_e \cdot f_l \cdot f_i \cdot f_c
  \label{eq:drake}
\end{equation}
donde $N_\star = 2\times10^{11}$ es el número de estrellas de la Vía
Láctea aptas para planetas con vida, $f_p = 1{,}0$ (fracción con
planetas, caso conservador que maximiza la población dado el producto
de filtros biológicos), $n_e = 0{,}22$ el número de planetas en zona
habitable por sistema \citep{fressin2013}, produciendo
$N_\text{hp} = N_\star \cdot f_p \cdot n_e = 4{,}4\times10^{10}$
planetas habitables potenciales, y $f_l$, $f_i$, $f_c$ son las
fracciones de vida, inteligencia y comunicación cuyo producto define el
escenario (Tabla \ref{tab:scenarios}).

Las posiciones galácticas se extraen del perfil de densidad continuo de
la ZGH \citep{lineweaver2004}:
\begin{equation}
  f(r) \propto \exp\!\left[-\frac{(r - r_\text{peak})^2}{2\sigma_r^2}\right],
  \quad r_\text{peak} = 25{,}000\,\text{al},\quad \sigma_r = 8{,}000\,\text{al}
  \label{eq:ghz}
\end{equation}
con una semialtura del disco de 300\,al muestreada de una distribución
exponencial con signo aleatorio.
Los tiempos de nacimiento se extraen uniformemente sobre la edad
estimada de la galaxia de $1{,}36\times10^{10}$\,años, y los tiempos
de vida tecnológicos siguen una distribución log-normal con mediana
$10^4$\,años \citep{sandberg2018}.

\begin{figure}
  \centering
  \includegraphics[width=\columnwidth]{output_analisis/analisis_ghz.png}
  \caption{Distribución radial de las civilizaciones generadas respecto al
           centro galáctico (seed\,=\,42, $N_c = 440$).
           Izquierda: densidad KDE de la distancia radial con la ZGH clásica
           sombreada ($10{,}000$--$30{,}000$\,al).
           Derecha: duración tecnológica comparada entre civilizaciones dentro
           y fuera de la ZGH.
           El $42{,}95\%$ de las civilizaciones se ubica dentro del anillo
           clásico de la ZGH, reflejando el perfil continuo de Lineweaver 2004
           (Ecuación~\ref{eq:ghz}) que distribuye nacimientos más allá de los
           límites convencionales.}
  \label{fig:ghz}
\end{figure}

\subsection{Heterogeneidad sociológica}
\label{sec:socio}

A cada civilización se le asigna uno de cuatro tipos sociológicos:
\emph{agresivo} (probabilidad 0{,}30), \emph{defensivo} (0{,}40),
\emph{pacifista} (0{,}20), \emph{curioso} (0{,}10).
En cada paso temporal, cada civilización adopta una de tres
estrategias activas -- \textsc{atacar}, \textsc{esconderse} o
\textsc{comunicar} -- muestreada de un vector de probabilidades
dependiente del tipo.
Las probabilidades base [\textsc{atacar}/\textsc{esconderse}/\textsc{comunicar}]
son: agresivo (0{,}65/0{,}10/0{,}25), defensivo (0{,}20/0{,}60/0{,}20),
pacifista (0{,}05/0{,}75/0{,}20) y curioso (0{,}10/0{,}15/0{,}75).
La estrategia modifica la detectabilidad de la civilización: como
\emph{presa}, los modificadores sobre $P_\text{det}$ son
$\times 1{,}20$ (\textsc{atacar}), $\times 0{,}001$
(\textsc{esconderse}) y $\times 3{,}00$ (\textsc{comunicar}); como
\emph{cazadora}, el modificador es $\times 1{,}20$ solo para
\textsc{atacar}.
La decisión de ataque final sigue una regla probabilística: la
probabilidad base de atacar es 0{,}90 (agresivo), 0{,}60 (defensivo),
0{,}10 (pacifista) y 0{,}30 (curioso), multiplicada por 1{,}50 si la
disparidad tecnológica relativa de la presa supera $\theta_\text{threat}$.

Las probabilidades base de estrategia se ajustan dinámicamente en
función de dos factores contextuales, ambos calculados a partir de
posiciones destruidas observadas localmente -- sin necesidad de
comunicación entre civilizaciones, preservando el silencio del Bosque
Oscuro:
\begin{enumerate}
  \item \emph{Presión ambiental inmediata} (radio de 500\,al):
        si hay $k$ civilizaciones destruidas dentro de 500\,al,
        la probabilidad de \textsc{esconderse} aumenta en
        $\delta = \min(0{,}40,\; 0{,}10\,k)$, mientras que
        \textsc{atacar} y \textsc{comunicar} se reducen cada una en
        $\delta/2$, con renormalización del vector resultante.
  \item \emph{Memoria histórica de zonas peligrosas} (radio de 2{,}000\,al):
        si existe al menos una civilización destruida dentro de un radio
        esférico de 2{,}000\,al, la probabilidad de \textsc{comunicar}
        se reduce en un 80 por ciento, y la masa redistribuida se
        transfiere íntegramente a \textsc{esconderse}, seguido de
        renormalización.
        El radio de memoria (2{,}000\,al) es mayor que el de presión
        inmediata (500\,al) para capturar el aprendizaje de amenazas
        en la vecindad galáctica amplia.
\end{enumerate}

\subsection{Probabilidad de detección}
\label{sec:detection}

La probabilidad de que la civilización $i$ detecte a la civilización
$j$ a distancia euclidiana $d$ en el tiempo $t$ es:
\begin{equation}
  P_\text{det}(d, \tau_i, \nu_i) =
    (1 - p_\text{fn})
    \cdot \frac{\rho(\tau_i,\nu_i)}{1 + (d/d_0)^2},
  \qquad d \leq c\,(t - t_{\text{nac},j})
  \label{eq:pdet}
\end{equation}
donde $\rho(\tau_i,\nu_i) = \min[1,\, r_\text{eff}(\tau_i,\nu_i)/r_\text{max}]$
es la capacidad tecnológica relativa del cazador, $d_0 = 100$\,al es
la escala de atenuación de la señal —correspondiente al alcance de
detección de fugas de radar con telescopios de próxima generación (SKA)
según \citet{loeb2007} y \citet{forgannichol2011}— y
$p_\text{fn} = 0{,}05$ la tasa de falsos negativos.
$P_\text{det} = 0$ si la condición causal no se satisface: ninguna
civilización es detectable antes de que su primera señal, viajando a
velocidad $c$, alcance al observador.
El perfil lorentziano garantiza $P_\text{det} \to 0$ cuando
$d \to \infty$.
La tasa de falsos positivos $p_\text{fp}$ no participa de la
Ecuación~\ref{eq:pdet}: un falso positivo corresponde, por definición,
a la detección de una señal donde no existe emisor, por lo que no puede
producir la localización de una civilización real.
La versión previa del modelo (v5) incluía $p_\text{fp}$ como término
aditivo, lo que imponía $P_\text{det} \to p_\text{fp} = 0{,}01$ cuando
$d \to \infty$ y convertía el ruido del detector en el canal dominante
de destrucción; el Apéndice~\ref{app:correction} documenta la
corrección y su impacto cuantitativo.
El radio efectivo de detección crece logísticamente:
\begin{equation}
  r_\text{eff}(\tau, \nu) =
    \frac{r_\text{max}}{1 + (r_\text{max}-1)\,e^{-\alpha\tau}}
    \cdot (1 + \max(0,\nu))
  \label{eq:reff}
\end{equation}
con $r_\text{max} = 1000$\,al \citep{troitskij1989,loeb2007} y tasa
de crecimiento $\alpha = 0{,}001$\,año$^{-1}$.
Con este valor, el tiempo de saturación al 50 por ciento de
$r_\text{max}$ es
$\tau_{1/2} = \ln(r_\text{max}-1)/\alpha \approx 6{,}900$\,años,
comparable a la mediana del tiempo de vida tecnológico ($L = 10^4$\,años).
Las civilizaciones con $\tau \ll \tau_{1/2}$ operan con $r_\text{eff}$
significativamente menor que $r_\text{max}$, lo que hace al parámetro
$\alpha$ relevante para civilizaciones jóvenes o de corta duración.
Un índice espacial cKDTree 3D \citep{maneewongvatana1999} limita las
consultas de vecindad a civilizaciones dentro del horizonte de viaje
de luz, reduciendo la complejidad por evento de $O(N^2)$ a
$O(N \log N)$.

\subsection{Tiempo de vida extendido}
\label{sec:lifetime}

El modelo distingue dos escalas temporales que controlan aspectos
diferentes de la dinámica.
El \emph{tiempo de vida tecnológico} $L \sim 10^4$\,años
(Sección~\ref{sec:civgen}) determina la ventana durante la cual una
civilización evoluciona y emite señales detectables; es el parámetro $L$
de la ecuación de Drake y controla la coexistencia temporal en sentido
estricto (Sección~\ref{sec:civgen}).
La \emph{ventana extendida} $L_\text{ext} = 1{,}43\times10^8$\,años
\citep{sandberg2018} determina la persistencia de la firma tecnológica
una vez que la civilización ha emergido, permitiendo que civilizaciones
cuyos períodos tecnológicos no se solapan puedan interactuar si sus
ventanas extendidas sí lo hacen.
Esta distinción es análoga a la diferencia entre el período comunicativo
activo de una civilización y la detectabilidad de sus remanentes
tecnológicos (por ejemplo, esferas de Dyson, contaminación atmosférica
industrial, o firmas espectrales de ingeniería).
En la práctica, el paso de dinámica del Bosque Oscuro
(Sección~\ref{sec:detection}) utiliza $L_\text{ext}$ como el tiempo
durante el cual una civilización permanece detectable como blanco,
mientras que las duraciones lognormales de la Sección~\ref{sec:civgen}
se emplean para el análisis de coexistencia temporal pre-dinámica.
Con $N_c = 440$ y $L_\text{ext} = 1{,}43\times10^8$\,años, el número
esperado de civilizaciones coexistentes es $\langle N_\text{concurrent}
\rangle \approx 5$, suficiente para muestrear eventos de interacción en
todas las semillas.
El valor de $L_\text{ext}$ se obtiene de la condición
$\langle N_\text{concurrent} \rangle \geq 5$ para $N_c = 440$ sobre una
ventana temporal $T_\text{gal} = 1{,}26\times10^{10}$\,años:
$L_\text{ext} \geq 5 \times T_\text{gal} / N_c \approx 1{,}43\times10^8$
\,años.
Este umbral garantiza al menos algunos pares interactuantes por semilla;
valores de $L_\text{ext}$ en el rango $10^7$--$10^9$\,años son
consistentes con las escalas de persistencia de firmas tecnológicas
discutidas por \citet{sandberg2018}.
$L_\text{ext}$ debe interpretarse como un parámetro libre calibrado
para viabilizar la simulación, no como una predicción derivada de la
literatura.

\subsection{Análisis de sensibilidad}
\label{sec:sensibilidad}

El análisis de sensibilidad de esta versión se centra en el parámetro
que la corrección v5\,$\to$\,v6 reveló como dominante: la tasa de
falsos positivos $p_\text{fp}$.
Se compara la dinámica completa (ambos regímenes, controles apareados)
bajo tres configuraciones: (a) el modelo v5 con el término
$p_\text{fp}$ aditivo; (b) el modelo sin piso ($p_\text{fp}$ excluido
del canal letal); y (c) el modelo final, que además impone la condición
causal de la Ecuación~\ref{eq:pdet}.
Los resultados se reportan en la Sección~\ref{sec:sens_results}.
La versión v5 incluía un análisis de Sobol global sobre cinco
parámetros ($\theta_\text{threat}$, $\alpha$, $f_\text{agresivo}$,
$r_\text{max}$, $d_0$; $N = 4{,}096$ muestras de Saltelli,
\citealt{saltelli2010}) que reportaba índices de primer orden
compatibles con cero para todos ellos.
Ese resultado era correcto pero engañoso: $p_\text{fp}$ —el parámetro
individual efectivamente dominante— no estaba entre los factores
variados, y el estimador (una réplica de 100 civilizaciones por punto
de muestreo) presentaba un coeficiente de variación de ruido Monte
Carlo cercano al 80 por ciento.
Un análisis de Sobol del modelo corregido, con réplicas promediadas
por punto y $p_\text{fn}$ incluido entre los factores, se deja
declarado como trabajo futuro.

\subsection{Simulaciones de control sin dinámica DF}
\label{sec:control}

Para aislar el efecto de la dinámica del Bosque Oscuro de los efectos
geométricos preexistentes de la ZGH, ejecutamos simulaciones de control
en las que la probabilidad de ataque se fuerza a cero para todos los
tipos sociológicos ($P_\text{ataque} = 0$).
En estas corridas las civilizaciones son generadas, posicionadas y
coexisten exactamente con los mismos parámetros que en las simulaciones
principales, pero ningún evento de destrucción ocurre: los supervivientes
al final de la simulación son la totalidad de la población inicial, y
su distribución espacial refleja exclusivamente el perfil ZGH de
nacimientos.

El cociente de control $r_\text{obs/P}^\text{ctrl}$ cuantifica entonces
la concentración espacial producida únicamente por la geometría ZGH.
La diferencia
\begin{equation}
  \Delta r = r_\text{obs/P}^\text{DF} - r_\text{obs/P}^\text{ctrl}
  \label{eq:delta_r}
\end{equation}
mide el efecto espacial \emph{neto} de la dinámica del Bosque Oscuro,
una vez eliminada la contribución geométrica de la ZGH.
Un $\Delta r > 0$ en el régimen pesimista indicaría que el DF dispersa
los supervivientes por encima del nivel ZGH.
Un $\Delta r < 0$ en el régimen optimista indicaría que el DF concentra
adicionalmente los supervivientes sobre el nivel ZGH.
Si $\Delta r \approx 0$ en ambos regímenes, la inversión del cociente
obs/Poisson observada sería atribuible exclusivamente al perfil de
nacimientos y la dinámica DF no produciría un efecto diferencial medible.
Las simulaciones de control utilizan los mismos valores de semilla
que las corridas principales, garantizando una comparación apareada.

\subsection{Optimización computacional}
\label{sec:optimization}

El paso de dinámica del Bosque Oscuro (Sección~\ref{sec:detection})
originalmente evaluaba la probabilidad de detección para cada par
cazador--presa de forma escalar dentro de un bucle Python anidado,
resultando en $\sim 10^5$ llamadas a \texttt{prob\_deteccion()} por
evento para $N_c = 10{,}000$.
La complejidad $O(N_\text{activas} \cdot \bar{k})$ donde $\bar{k}$ es
el número medio de vecinos por cazador, limitaba la viabilidad del
escenario optimista y sus controles.

Se implementaron dos optimizaciones complementarias:
\begin{enumerate}
  \item \emph{Vectorización del bucle interno de presas.}
        Para cada cazador, las $n$ presas candidatas se procesan
        simultáneamente: las distancias 3D se calculan con operaciones
        vectoriales de \textsc{numpy}, la probabilidad de detección se
        evalúa en batch mediante \texttt{prob\_deteccion\_batch(dists,
        edad, nivel, p)} —una versión vectorizada de la
        Ecuación~\ref{eq:pdet} que opera sobre arrays de $n$ distancias—,
        y las $n$ pruebas de Bernoulli se ejecutan con una sola llamada a
        \texttt{np.random.random(n)}.
        Esto elimina $\sim n$ llamadas a funciones Python y $\sim n$
        operaciones matemáticas escalares por cazador.
  \item \emph{Extracción temprana de arrays numpy.}
        Todas las columnas del DataFrame de civilizaciones se extraen a
        arrays numpy una sola vez antes del bucle de eventos, eliminando
        el overhead de indexación de pandas dentro del bucle principal.
\end{enumerate}
La combinación de ambas optimizaciones produce un speedup de
$\sim 5\times$ respecto de la implementación escalar, haciendo viable
el control optimista.
En el modelo corregido, donde la mayoría de las civilizaciones
sobrevive y permanece activa durante toda la simulación, una corrida
con $N_c = 10{,}000$ toma $\sim 6$\,minutos en un núcleo de CPU.

% ─────────────────────────────────────────────────────────────────────────────
\section{Resultados}
\label{sec:results}
% ─────────────────────────────────────────────────────────────────────────────

\subsection{Tasas de destrucción}
\label{sec:destruction}

La Tabla \ref{tab:results} resume las estadísticas de destrucción para
ambos escenarios (50 semillas pesimista, 10 semillas optimista).

\begin{table*}
  \centering
  \caption{Estadísticas resumen para los escenarios pesimista
           (50 semillas DF + 50 de control apareadas) y optimista
           (10 semillas DF + 3 de control).
           Las incertidumbres son $\pm 1\sigma$ entre semillas.
           $\tilde{d}_\text{kill}$ es la mediana de la distancia
           cazador--presa en los eventos de destrucción.
           $\Delta r = r_\text{obs/P}^\text{DF} - r_\text{obs/P}^\text{ctrl}$
           cuantifica el efecto espacial neto del Bosque Oscuro.}
  \label{tab:results}
  \begin{tabular}{lccccccc}
    \toprule
    Escenario & $N_c$ & Destruidas & Tasa (\%) &
    $\tilde{d}_\text{kill}$ (al) &
    $\bar{d}_\text{nn}$ (al) & $r_\text{obs/P}$ & $\Delta r$ \\
    \midrule
    Pesimista & 440 & $2{,}6 \pm 1{,}5$ & $0{,}59 \pm 0{,}34$ &
    $736$ & $1733 \pm 53$ & $1{,}417 \pm 0{,}043$ &
    $+0{,}0045 \pm 0{,}0032^\ddagger$ \\
    Optimista  & $10{,}000^\dagger$ & $2287 \pm 46$ & $22{,}9 \pm 0{,}5$ &
    $2408$ & $553 \pm 4$  & $1{,}177 \pm 0{,}008$ &
    $+0{,}060 \pm 0{,}008^{\dagger\dagger}$ \\
    \bottomrule
  \end{tabular}
  \begin{flushleft}
    \footnotesize
    $^\dagger$ Estimación Drake cruda: $8{,}8\times10^6$; acotada en $10^4$
    por viabilidad computacional (véase Sección \ref{sec:limitations}).\\
    $^\ddagger$ Comparación apareada por semilla ($n = 50$;
    error estándar de la media: $0{,}0005$): efecto de dispersión
    positivo pero minúsculo. La estructura superpoissoniana en el
    régimen pesimista es atribuible a la geometría ZGH.\\
    $^{\dagger\dagger}$ $\Delta r_\text{opt} = +0{,}060 \pm 0{,}008$
    ($7{,}8\sigma$; control: $n = 3$ semillas,
    $r_\text{obs/P}^\text{ctrl} = 1{,}117 \pm 0{,}001$).
    El Bosque Oscuro dispersa a los supervivientes por encima del nivel
    geométrico ZGH en el régimen de alta densidad.
  \end{flushleft}
\end{table*}

Bajo parámetros pesimistas (50 semillas) la tasa de destrucción es
$0{,}59 \pm 0{,}34$ por ciento: con $\langle N_\text{concurrent}
\rangle \approx 5$ civilizaciones coexistentes y detección acotada por
la Ecuación~\ref{eq:pdet}, los encuentros letales son raros.
La mediana de la distancia de destrucción (736\,al) es consistente con
la escala física del modelo ($d_0 = 100$\,al, $r_\text{max} =
1000$\,al).
Bajo parámetros optimistas (10 semillas) la tasa asciende a
$22{,}9 \pm 0{,}5$ por ciento, un aumento de un factor $\sim 40$ que
refleja la tasa de encuentro dramáticamente mayor con
$N_c = 10{,}000$ ($\langle N_\text{concurrent} \rangle \approx 113$).
La mediana de la distancia de destrucción crece a $2{,}408$\,al: en el
régimen denso, la acumulación de un gran número de pares sobre la cola
lorentziana de la Ecuación~\ref{eq:pdet} hace que destrucciones a
varios miles de años luz, individualmente improbables, se vuelvan
colectivamente frecuentes.

\begin{figure}
  \centering
  \includegraphics[width=\columnwidth]{output_analisis/optimista/analisis_destrucciones.png}
  \caption{Distribución por estado post-paso5 (izquierda) y perfil radial
           de civilizaciones activas y destruidas (derecha) para la semilla
           de referencia (seed\,=\,42, escenario \emph{optimista},
           $N_c = 10{,}000$, $2{,}208$ destruidas).
           La distancia media de las destruidas al centro galáctico
           ($14{,}386$\,al) es menor que la de las activas
           ($21{,}757$\,al): la destrucción es más probable hacia el
           interior galáctico, donde la densidad tridimensional de
           civilizaciones (anillo ZGH interno más bulbo) produce mayores
           tasas de encuentro.}
  \label{fig:destrucciones}
\end{figure}

\subsection{Agrupamiento espacial y el cociente obs/Poisson}
\label{sec:clustering}

Para cada semilla y escenario calculamos la distancia media al vecino
más cercano $\bar{d}_\text{nn}$ entre civilizaciones activas
supervivientes y la comparamos con la expectativa de Poisson para un
campo aleatorio homogéneo de la misma densidad:
\begin{equation}
  d_\text{Poisson} = 0{,}5538 \left(\frac{V_\text{eff}}{N_\text{activas}}\right)^{1/3}
  \label{eq:poisson_nn}
\end{equation}
donde $V_\text{eff}$ es el volumen galáctico efectivo muestreado.
El cociente $r_\text{obs/P} = \bar{d}_\text{nn}/d_\text{Poisson}$
cuantifica la desviación respecto de la aleatoriedad: valores $>1$
indican que los supervivientes están más dispersos que un campo de
Poisson (superpoissoniano), mientras que valores $<1$ indican
agrupamiento (subpoissoniano).
Para aislar la contribución dinámica del DF de la concentración
geométrica de la ZGH, los resultados de las simulaciones principales
se contrastan con las simulaciones de control descritas en la
Sección~\ref{sec:control} y el efecto diferencial $\Delta r$
(ecuación~\ref{eq:delta_r}).

\begin{figure}
  \centering
  \includegraphics[width=\columnwidth]{output_analisis/distribucion_vecino_mas_cercano.png}
  \caption{Distribución de la distancia al vecino más cercano en 3D para
           las civilizaciones activas supervivientes (seed\,=\,42, escenario
           pesimista, $n = 439$).
           La línea roja marca la media observada ($1792$\,al), la naranja
           la mediana ($1517$\,al), y la verde la distancia esperada bajo un
           proceso de Poisson homogéneo ($1222$\,al).
           El cociente $r_\text{obs/P} = 1{,}47$ refleja una desviación
           superpoissoniana atribuible principalmente al perfil continuo
           de la ZGH (véase el control apareado,
           Sección~\ref{sec:control_results}).}
  \label{fig:vecino}
\end{figure}

\subsubsection{Distribución ZGH de referencia (control sin DF)}
\label{sec:control_results}

Las simulaciones de control ($P_\text{ataque} = 0$, 50 semillas
apareadas) producen
$r_\text{obs/P}^\text{ctrl} = 1{,}413 \pm 0{,}043$.
El efecto diferencial neto del Bosque Oscuro, calculado como diferencia
apareada por semilla, es
$\Delta r = r_\text{obs/P}^\text{DF} - r_\text{obs/P}^\text{ctrl}
= +0{,}0045 \pm 0{,}0032$ (error estándar de la media: $0{,}0005$,
$n = 50$ pares).
El apareamiento por semilla —catálogos idénticos con y sin dinámica de
destrucción— cancela la varianza geométrica entre realizaciones y
permite resolver efectos dos órdenes de magnitud menores que la
dispersión entre semillas ($\sigma_\text{ctrl} = 0{,}043$).
El efecto detectado es positivo (dispersión) pero minúsculo: la
eliminación del $0{,}59\%$ de las civilizaciones apenas perturba la
estructura espacial, y la señal superpoissoniana observada
($r_\text{obs/P} > 1$) refleja esencialmente la concentración
geométrica de los nacimientos en el anillo de habitabilidad galáctica.

\subsubsection{Régimen pesimista: supervivientes superpoissonianos}

En el escenario pesimista, $r_\text{obs/P} = 1{,}417 \pm 0{,}043$
en 50 semillas (calculado sobre las civilizaciones activas
supervivientes en cada semilla).
Dado que $\Delta r \approx 0$, esta estructura superpoissoniana es
atribuible casi enteramente al perfil continuo de la ZGH
(Ecuación~\ref{eq:ghz}), que concentra los nacimientos en un anillo
gaussiano alrededor de $r_\text{peak} = 25{,}000$\,al.
Esto es cualitativamente consistente con el modelo de supervivencia
basado en vacíos de \citet{forgan2011}, aunque nuestro análisis de
control muestra que dichos vacíos reflejan la geometría de
nacimientos, no la dinámica de destrucción.

\subsubsection{Régimen optimista: dispersión neta sobre el nivel ZGH}
\label{sec:inversion}

En el escenario optimista, $r_\text{obs/P} = 1{,}177 \pm 0{,}008$ en
las 10 semillas independientes.
El control en este régimen ($n = 3$ semillas, $N_c = 10{,}000$,
$P_\text{ataque} = 0$) produce
$r_\text{obs/P}^\text{ctrl} = 1{,}117 \pm 0{,}001$: el efecto
diferencial es
$\Delta r_\text{opt} = r_\text{obs/P}^\text{DF} -
r_\text{obs/P}^\text{ctrl} = +0{,}060 \pm 0{,}008$ ($7{,}8\sigma$
usando la dispersión entre semillas como incertidumbre).
El Bosque Oscuro \emph{dispersa} a los supervivientes por encima del
nivel geométrico ZGH.
El mecanismo es directo: la probabilidad de detección decae con la
distancia (Ecuación~\ref{eq:pdet}), por lo que los eventos de
destrucción eliminan preferentemente pares cercanos; cada destrucción
tiende a remover al vecino más próximo de alguna civilización,
desplazando la distribución de $d_\text{nn}$ de los supervivientes
hacia valores mayores.
El efecto crece monotónicamente con la densidad
($+0{,}0045 \to +0{,}060$ entre regímenes, un factor $\sim 13$)
porque la fracción destruida crece con la tasa de encuentros
($0{,}59\% \to 22{,}9\%$).

Este resultado invierte el hallazgo de la versión previa del modelo,
que reportaba concentración masiva
($\Delta r_\text{opt} = -0{,}538 \pm 0{,}020$) con supervivientes
subpoissonianos ($r_\text{obs/P} = 0{,}578$).
La Sección~\ref{sec:sens_results} demuestra que aquella señal era un
artefacto del piso de falsos positivos: con detección efectiva a
cualquier distancia, la supervivencia dejaba de depender de la
proximidad entre pares y pasaba a estar dominada por la estrategia de
ocultamiento, cuya adopción se correlaciona espacialmente con las
zonas de destrucción previa, produciendo el agrupamiento espurio.

\citet{foucher2023} reportó distribuciones subpoissonianas usando el
índice de dispersión $R = \sigma^2/\mu$ y la función de correlación de
pares $g(r)$, atribuyéndolas a dinámica depredador--presa.
Nuestro modelo corregido no reproduce ese régimen: la destrucción por
proximidad produce el efecto contrario (dispersión).
La comparación sugiere cautela: señales de agrupamiento en modelos de
tipo depredador--presa pueden depender críticamente de la estructura
del canal de detección a larga distancia, como ilustra la corrección
v5\,$\to$\,v6 de este trabajo.

\subsection{Sensibilidad: el papel dominante de $p_\text{fp}$}
\label{sec:sens_results}

La Tabla~\ref{tab:sens_pfp} compara la dinámica completa bajo las tres
configuraciones del canal de detección definidas en la
Sección~\ref{sec:sensibilidad}.

\begin{table}
  \centering
  \caption{Análisis de sensibilidad de un factor sobre el término de
           falsos positivos. Cada fila usa las mismas semillas
           (50 pesimista, 10 optimista) y los mismos controles.
           $\tilde{d}_\text{kill}$: mediana de la distancia de
           destrucción (régimen indicado).}
  \label{tab:sens_pfp}
  \begin{tabular}{lccc}
    \toprule
    Configuración & Tasa pes. (\%) & Tasa opt. (\%) & $\Delta r_\text{opt}$ \\
    \midrule
    v5 ($p_\text{fp}$ aditivo)   & $16{,}1 \pm 1{,}6$ & $85{,}8 \pm 0{,}2$ & $-0{,}538 \pm 0{,}020$ \\
    sin piso $p_\text{fp}$        & $0{,}69 \pm 0{,}38$ & $23{,}0 \pm 0{,}3$ & $+0{,}061 \pm 0{,}006$ \\
    final (+ causalidad)          & $0{,}59 \pm 0{,}34$ & $22{,}9 \pm 0{,}5$ & $+0{,}060 \pm 0{,}008$ \\
    \bottomrule
  \end{tabular}
  \begin{flushleft}
    \footnotesize
    Mediana de la distancia de destrucción en el régimen pesimista:
    $\sim 32{,}000$\,al (v5) frente a $736$\,al (final); en el régimen
    optimista: $2{,}408$\,al (final).
    En v5, el 30\% de las destrucciones violaba la condición causal
    $d \leq c\,(t - t_\text{nac})$.
  \end{flushleft}
\end{table}

El término aditivo de falsos positivos controlaba la casi totalidad de
la dinámica del modelo v5: su eliminación reduce la tasa de destrucción
pesimista en un factor $\sim 27$ y la optimista en un factor
$\sim 3{,}7$, e invierte el signo del efecto espacial diferencial.
El mecanismo es elemental: con
$P_\text{det} \to p_\text{fp} = 0{,}01$ cuando $d \to \infty$, cada par
cazador--presa activo constituía un ensayo de Bernoulli con
probabilidad $\geq 0{,}01$ por evento, independiente de la distancia;
a $d = 2{,}000$\,al la contribución lorentziana ($\approx 0{,}0024$) ya
era cuatro veces menor que el piso.
La consecuencia observable era una distribución de distancias de
destrucción físicamente incoherente (mediana $\sim 32{,}000$\,al, con
el 100\% de los eventos más allá de $2{,}000$\,al), que ninguna de las
métricas agregadas reportadas en v5 capturaba.
El análisis de Sobol de v5, que no incluía $p_\text{fp}$ entre sus
cinco factores, reportaba correctamente $S_1 \approx 0$ para todos los
parámetros variados; la conclusión de que ``ningún parámetro individual
domina'' era válida solo dentro del subespacio muestreado.
La condición causal, añadida en la configuración final, tiene un efecto
menor una vez eliminado el piso (reducción adicional de $\sim 15$ por
ciento en la tasa pesimista), pero elimina una clase entera de eventos
no físicos.
Recomendamos que los modelos Monte Carlo de la paradoja de Fermi
reporten la distribución de distancias de interacción como diagnóstico
estándar de coherencia física.

% ─────────────────────────────────────────────────────────────────────────────
\section{Discusión}
\label{sec:discussion}
% ─────────────────────────────────────────────────────────────────────────────

\subsection{Dispersión modesta y monotónica con la densidad}

El hallazgo central es que el efecto espacial neto de la dinámica del
Bosque Oscuro es una dispersión de los supervivientes, pequeña en
términos absolutos y creciente con la densidad poblacional.
En ambos regímenes la estructura espacial está dominada por la
geometría de nacimientos de la ZGH: los cocientes obs/Poisson de las
corridas con dinámica activa ($1{,}417$ y $1{,}177$) difieren de sus
controles ($1{,}413$ y $1{,}117$) en menos de un 6 por ciento.
Sobre ese fondo geométrico, el contraste apareado resuelve una señal
de dispersión sistemática: $\Delta r = +0{,}0045 \pm 0{,}0032$ en baja
densidad y $\Delta r = +0{,}060 \pm 0{,}008$ en alta densidad.
El mecanismo es la localidad de la destrucción: como
$P_\text{det}$ decae con la distancia, los eventos letales eliminan
preferentemente pares próximos, y cada eliminación desplaza la
distribución de distancias al vecino más cercano hacia valores
mayores.
La magnitud del efecto sigue a la fracción destruida
($0{,}59\% \to 22{,}9\%$), que a su vez sigue a la tasa de encuentros:
en el escenario pesimista la separación media entre civilizaciones
coexistentes ($\sim 2{,}200$\,al) supera $r_\text{max}$, mientras que
en el optimista ($\sim 550$\,al) queda por debajo, y la cola
lorentziana acumulada sobre $\sim 113$ civilizaciones concurrentes
produce destrucciones frecuentes hasta la escala de unos pocos miles
de años luz.

Ninguno de los dos regímenes muestra la concentración subpoissoniana
reportada por la versión previa del modelo.
Aquel resultado ($\Delta r = -0{,}54$, supervivientes a
$r_\text{obs/P} = 0{,}58$) dependía por completo del piso de falsos
positivos (Sección~\ref{sec:sens_results}): al hacer la detección
efectiva a cualquier distancia, el piso desacoplaba la supervivencia
de la proximidad y la acoplaba a la estrategia de ocultamiento,
generando un agrupamiento que la corrección elimina.

\subsection{Comparación con trabajos previos}

\citet{forgan2009} encontró $\sim 361$ civilizaciones bajo parámetros
de ``Vida Rara'' con $P_\text{destroy} = 0{,}5$, comparable en escala
a nuestro escenario pesimista.
Su modelo no calculó estadísticas espaciales de los supervivientes.
El presente modelo extiende su marco de trabajo con una función de
detección físicamente motivada, heterogeneidad sociológica y análisis
espacial explícito.

\citet{foucher2023} reportó supervivientes subpoissonianos en modelos
ecológicos de alta densidad.
Nuestro modelo corregido produce el efecto contrario (dispersión por
eliminación de vecinos próximos); la discrepancia sugiere que el signo
de la señal espacial en esta clase de modelos depende críticamente del
comportamiento del canal de detección a larga distancia, y refuerza la
recomendación de reportar la distribución de distancias de interacción
(Sección~\ref{sec:sens_results}).

\subsection{El parámetro libre del umbral de ataque}

Ningún estudio publicado restringe $\theta_\text{threat}$
empíricamente \citep{devladar2013}.
En el modelo corregido, $\theta_\text{threat}$ modula la probabilidad
de ataque condicionada a una detección previa; dado que la detección
está fuertemente acotada por la distancia (Ecuación~\ref{eq:pdet}),
su efecto sobre la estructura espacial es de segundo orden frente a la
geometría de encuentros.
Una cuantificación formal (Sobol sobre el modelo corregido, con
réplicas promediadas) se deja como trabajo futuro.
La inferencia bayesiana contra el registro de no-detección del SETI
podría proporcionar priores informativos y constituye una extensión
natural.

% ─────────────────────────────────────────────────────────────────────────────
\section{Limitaciones y trabajo futuro}
\label{sec:limitations}
% ─────────────────────────────────────────────────────────────────────────────

\textit{Límite computacional.}
La estimación cruda de Drake para el escenario optimista es
$N_c = 8{,}8\times10^6$, acotada en $10^4$ (0{,}11 por ciento) por
viabilidad computacional.
Las optimizaciones descritas en la Sección~\ref{sec:optimization}
(vectorización del bucle interno, extracción temprana de arrays)
redujeron el tiempo de ejecución en un factor $\sim 5\times$,
viabilizando el control optimista.
Escalar a $N_c \sim 10^5$--$10^6$ requeriría vectorizar el bucle
externo sobre cazadores (actualmente en Python), acotar el radio de
consulta espacial al alcance físico de detección y paralelizar sobre
múltiples núcleos, lo que se deja para trabajo futuro.
El límite actual afecta las tasas absolutas pero no la tendencia
cualitativa —dispersión creciente con la densidad—, que depende del
cociente entre la separación media y $r_\text{max}$.

\textit{Movimiento propio estelar.}
Las posiciones fijas ignoran la deriva estelar
\citep[$\sim 30$\,km\,s$^{-1}$;][]{carrollnellenback2019},
que difundiría la concentración de la ZGH y reduciría la magnitud de
la señal subpoissoniana.

\textit{Parámetros libres.}
La inferencia bayesiana contra observaciones SETI reemplazaría los
valores ad hoc de $\theta_\text{threat}$, las fracciones sociológicas
y las tasas de ruido de detección con posteriores basados en datos.

\textit{Escalas de detección y memoria histórica.}
La escala de atenuación $d_0 = 100$\,al y el radio de memoria histórica
$R_\text{mem} = 2{,}000$\,al difieren en un factor $20\times$.
A $d = 2{,}000$\,al, la atenuación lorentziana es $1/401 \approx
0{,}0025$: el mecanismo de memoria histórica —que reduce la
comunicación en regiones donde se han observado destrucciones— opera
mayormente a distancias donde la detectabilidad es baja, aunque en el
régimen denso la mediana de las destrucciones ($2{,}408$\,al) coincide
con la escala de memoria, por lo que el mecanismo actúa sobre la
población relevante.
El efecto de silencio inducido por la memoria es débil en términos
absolutos, aunque puede ser relevante para civilizaciones
anómalamente detectables (estrategia \textsc{comunicar}, $\times 3{,}0$)
o con niveles tecnológicos excepcionalmente altos.
En el modelo v5, la verificación de robustez con memoria desactivada
($R_\text{mem} \to 0$, semilla de referencia) varió la tasa de
destrucción en $< 2$ puntos porcentuales y $\Delta r$ en $< 0{,}01$;
repetir esta verificación sobre el modelo corregido queda como trabajo
futuro.

\textit{Geometría del disco galáctico.}
La escala de altura del disco utilizada (300\,al) corresponde al disco
delgado de gas y polvo; el disco estelar grueso ($\sim 1000$\,al)
aumentaría las distancias 3D típicas en un factor $\sim 1{,}5$ y
reduciría las tasas de interacción.
Este efecto es degenerado con una reducción efectiva de $r_\text{max}$
y no altera la inversión cualitativa del cociente obs/Poisson, pero
los valores absolutos de las distancias al vecino más cercano deben
interpretarse como cotas inferiores.
Trabajo futuro debería incorporar un modelo de doble disco (thin + thick)
con escalas ajustadas a la población estelar de la Vía Láctea.

\textit{Predicción observacional.}
El crecimiento monotónico de la dispersión con la densidad
($\Delta r: +0{,}005 \to +0{,}060$) constituye, en principio, una
predicción contrastable: si la HBO opera en galaxias espirales, las
fuentes tecnológicas supervivientes en regiones de alta densidad
estelar deberían exhibir un exceso de regularidad espacial —distancias
al vecino más cercano sistemáticamente mayores que las de la población
de nacimientos subyacente— del orden del 5 por ciento sobre el nivel
geométrico.
Debemos ser francos sobre la dificultad: una señal de esta magnitud es
mucho más exigente observacionalmente que la concentración masiva que
predecía el modelo v5, requiere estimar la geometría de nacimientos de
forma independiente, y probablemente quede fuera del alcance de los
catálogos de próxima generación (SKA, WISE) salvo para poblaciones
muy numerosas de candidatos.
La predicción robusta del modelo corregido es, más bien, negativa: la
HBO \emph{no} genera firmas espaciales fuertes; el silencio galáctico
que produce es estadísticamente casi indistinguible de la geometría
de habitabilidad.

% ─────────────────────────────────────────────────────────────────────────────
\section{Conclusiones}
\label{sec:conclusions}
% ─────────────────────────────────────────────────────────────────────────────

\begin{enumerate}
  \item Bajo parámetros de Drake pesimistas ($N_c = 440$, 50 semillas),
        el mecanismo del Bosque Oscuro con detección físicamente
        acotada destruye solo $0{,}59 \pm 0{,}34$\,\% de las
        civilizaciones, con distancias de destrucción coherentes con la
        escala del modelo (mediana $736$\,al).
        La distribución de supervivientes es superpoissoniana
        ($r_\text{obs/P} = 1{,}417 \pm 0{,}043$), pero el contraste
        apareado con 50 controles ($P_\text{ataque} = 0$) muestra que
        esta estructura es la de los nacimientos ZGH: el efecto neto de
        la dinámica es $\Delta r = +0{,}0045 \pm 0{,}0032$, positivo
        pero minúsculo.

  \item Bajo parámetros optimistas ($N_c = 10{,}000$, 10 semillas DF +
        3 de control), la tasa de destrucción asciende a
        $22{,}9 \pm 0{,}5$\,\% y el efecto diferencial es
        $\Delta r_\text{opt} = +0{,}060 \pm 0{,}008$ ($7{,}8\sigma$):
        la dinámica DF \emph{dispersa} a los supervivientes por encima
        del nivel geométrico ZGH, porque las destrucciones son locales
        y eliminan preferentemente al vecino más cercano.

  \item El efecto espacial neto del Bosque Oscuro es, por tanto, una
        dispersión modesta y monotónica con la densidad
        ($+0{,}005 \to +0{,}060$), no la concentración masiva
        ($\Delta r = -0{,}54$) reportada por la versión previa del
        modelo.
        Aquella señal era un artefacto del término aditivo de falsos
        positivos, que imponía una probabilidad de detección
        $\geq 0{,}01$ a cualquier distancia y generaba más del 96 por
        ciento de las destrucciones a distancias no físicas
        (mediana $\sim 32{,}000$\,al), con un 30 por ciento de eventos
        acausales.

  \item El análisis de sensibilidad de un factor demuestra que
        $p_\text{fp}$ era el parámetro individual dominante del modelo
        v5 —su eliminación cambia las tasas de destrucción en factores
        de $4$--$27$ e invierte el signo de $\Delta r$—, pese a que el
        análisis de Sobol de esa versión (que no lo incluía entre sus
        factores) reportaba que ningún parámetro dominaba.
        La lección es general: la coherencia física de los términos de
        ruido de detección debe verificarse directamente, y la
        distribución de distancias de interacción es el diagnóstico
        natural.

  \item Cuatro parámetros del modelo siguen sin restricción
        observacional directa ($\theta_\text{threat}$, $d_0$, las
        fracciones sociológicas, $p_\text{fn}$); las tasas absolutas
        deben interpretarse como exploraciones del espacio de fases.
        La conclusión cualitativa robusta es negativa y de relevancia
        directa para SETI: incluso cuando la dinámica del Bosque Oscuro
        opera de forma ubicua (régimen denso), su firma espacial es
        débil, y el silencio que produce es casi indistinguible de la
        geometría de habitabilidad galáctica.
\end{enumerate}

El código, los datos y el historial completo de la corrección
v5\,$\to$\,v6 están depositados en
\url{https://github.com/CienciaEstelar/dark-forest-mc}.

\section*{Agradecimientos}

El autor agradece a las comunidades de código abierto detrás de
\textsc{numpy}, \textsc{scipy}, \textsc{pandas}, \textsc{SALib}
y \textsc{pytest}.
Esta investigación no recibió financiamiento externo.

\section*{Disponibilidad de datos}

Los resultados definitivos del modelo corregido
(\texttt{docs/resultados\_definitivos\_fix\_pfp\_causal.json}: 50
semillas pesimistas, 50 controles apareados y 10 semillas optimistas,
con estadísticos por semilla), la trazabilidad completa de la
corrección (\texttt{docs/ROADMAP\_CORRECCION.md}, Revisión 6), la
suite de pruebas (regresión con semillas 42, 7, 13, 99, 256, más
pruebas de propiedades físicas del canal de detección:
ausencia de piso y causalidad) y el código fuente están depositados en
\url{https://github.com/CienciaEstelar/dark-forest-mc}.

\appendix
\section{Parámetros del modelo}
\label{app:params}

\begin{table}
  \centering
  \caption{Parámetros del filtro de Drake para ambos escenarios.}
  \label{tab:scenarios}
  \begin{tabular}{lccc}
    \toprule
    Escenario & $f_l$ & $f_i$ & $f_c$ \\
    \midrule
    Pesimista & $10^{-3}$ & $10^{-3}$ & $10^{-2}$ \\
    Optimista  & $0{,}10$    & $0{,}01$    & $0{,}20$    \\
    \bottomrule
  \end{tabular}
\end{table}

\begin{table}
  \centering
  \caption{Parámetros fijos del modelo.}
  \label{tab:params}
  \begin{tabular}{lll}
    \toprule
    Parámetro & Valor & Referencia \\
    \midrule
    $N_\star$ & $2\times10^{11}$ & -- \\
    $r_\text{max}$ & 1000\,al & \citet{troitskij1989,loeb2007} \\
    $r_\text{peak}$ (ZGH) & 25{,}000\,al & \citet{lineweaver2004} \\
    $\sigma_r$ (ZGH) & 8{,}000\,al & \citet{lineweaver2004} \\
    $L_\text{ext}$ & $1{,}43\times10^8$\,años & \citet{sandberg2018} \\
    $\theta_\text{threat}$ & 0{,}70 & libre \\
    $\alpha$ & 0{,}001\,año$^{-1}$ & libre \\
    $d_0$ (atenuación) & 100\,al & \citet{loeb2007}; \citet{forgannichol2011} \\
    $p_\text{fn}$ & 0{,}05 & libre \\
    $p_\text{fp}$ & 0{,}01$^{*}$ & libre \\
    \bottomrule
  \end{tabular}
  \begin{flushleft}
    \footnotesize
    $^{*}$ Desde la corrección v6, $p_\text{fp}$ no participa del canal
    de detección de blancos reales (Ecuación~\ref{eq:pdet} y
    Apéndice~\ref{app:correction}); se lista por trazabilidad con la
    versión v5.
  \end{flushleft}
\end{table}

\section{Historial de la corrección v5 $\to$ v6}
\label{app:correction}

La versión v5 de este modelo (DOI Zenodo:
\texttt{10.5281/zenodo.20451755}) incluía la tasa de falsos positivos
como término aditivo en la probabilidad de detección:
$P_\text{det}^\text{v5} = p_\text{fp} + (1 - p_\text{fn} -
p_\text{fp})\,\rho/(1 + (d/d_0)^2)$, con
$P_\text{det}^\text{v5} \to p_\text{fp} = 0{,}01$ cuando
$d \to \infty$.
Una auditoría posterior del código reveló tres consecuencias no
reportadas en v5:
(i) en el régimen pesimista, el 100 por ciento de las destrucciones
ocurría a $d > 2{,}000$\,al, con mediana $\sim 32{,}000$\,al —
comparable al radio galáctico y $32\times$ el radio máximo de
detección $r_\text{max} = 1000$\,al que el propio modelo justifica con
\citet{troitskij1989} y \citet{loeb2007};
(ii) el 30 por ciento de las destrucciones violaba la condición causal
$d \leq c\,(t - t_\text{nac})$, pues el modelo solo acotaba el
horizonte de información del cazador, no el tiempo de viaje de la
señal de la presa; y
(iii) la eliminación del término ($p_\text{fp} = 0$) colapsaba la tasa
de destrucción pesimista de $16{,}1$ a $0{,}7$ por ciento y la
optimista de $85{,}8$ a $23{,}0$ por ciento, invirtiendo el signo del
efecto espacial diferencial (Tabla~\ref{tab:sens_pfp}).
La corrección v6 (i) excluye $p_\text{fp}$ del canal de detección de
blancos reales —un falso positivo es la detección de una señal sin
emisor y no puede localizar una civilización existente— y (ii) impone
la condición causal en la dinámica.
Ambas propiedades quedan fijadas por pruebas automatizadas en el
repositorio público.
Todos los resultados de este artículo corresponden al modelo
corregido; los valores v5 se reportan únicamente en la
Tabla~\ref{tab:sens_pfp} y en este apéndice como registro de la
corrección.

\begin{thebibliography}{99}

\bibitem[\protect\citeauthoryear{Brin}{1983}]{brin1983}
Brin G.~D., 1983, QJRAS, 24, 283

\bibitem[\protect\citeauthoryear{Carroll-Nellenback et al.}{2019}]{carrollnellenback2019}
Carroll-Nellenback J., Frank A., Wright J., Scharf C., 2019, AJ, 158, 117.
\newblock \doi{10.3847/1538-3881/ab31a3}

\bibitem[\protect\citeauthoryear{Cixin}{2008}]{cixin2008}
Cixin L., 2008, The Dark Forest. Chongqing Publishing House

\bibitem[\protect\citeauthoryear{de Vladar}{2013}]{devladar2013}
de Vladar H.~P., 2013, IJA, 12, 53.
\newblock \doi{10.1017/S1473550412000407}

\bibitem[\protect\citeauthoryear{Fermi}{1950}]{fermi1950}
Fermi E., 1950, observaciones no publicadas, Los Álamos

\bibitem[\protect\citeauthoryear{Fressin et al.}{2013}]{fressin2013}
Fressin F. et al., 2013, ApJ, 766, 81.
\newblock \doi{10.1088/0004-637X/766/2/81}

\bibitem[\protect\citeauthoryear{Forgan}{2009}]{forgan2009}
Forgan D.~H., 2009, IJA, 8, 121.
\newblock \doi{10.1017/S1473550408004321}

\bibitem[\protect\citeauthoryear{Forgan}{2011}]{forgan2011}
Forgan D.~H., 2011, IJA, 10, 341.
\newblock \doi{10.1017/S1473550411000127}

\bibitem[\protect\citeauthoryear{Forgan \& Nichol}{2011}]{forgannichol2011}
Forgan D.~H., Nichol R.~C., 2011, IJA, 10, 77.
\newblock \doi{10.1017/S1473550410000236}

\bibitem[\protect\citeauthoryear{Foucher}{2023}]{foucher2023}
Foucher F., 2023, AcAau, 204, 785.
\newblock \doi{10.1016/j.actaastro.2022.12.007}

\bibitem[\protect\citeauthoryear{Hart}{1975}]{hart1975}
Hart M.~H., 1975, QJRAS, 16, 128

\bibitem[\protect\citeauthoryear{Lineweaver, Fenner \& Gibson}{2004}]{lineweaver2004}
Lineweaver C.~H., Fenner Y., Gibson B.~K., 2004, Science, 303, 59.
\newblock \doi{10.1126/science.1092322}

\bibitem[\protect\citeauthoryear{Loeb \& Zaldarriaga}{2007}]{loeb2007}
Loeb A., Zaldarriaga M., 2007, JCAP, 2007, 020.
\newblock \doi{10.1088/1475-7516/2007/01/020}

\bibitem[\protect\citeauthoryear{Maneewongvatana \& Mount}{1999}]{maneewongvatana1999}
Maneewongvatana S., Mount D.~M., 1999, preimpresión (arXiv:cs/9901013)

\bibitem[\protect\citeauthoryear{Saltelli et al.}{2010}]{saltelli2010}
Saltelli A. et al., 2010, Comput.~Phys.~Commun., 181, 259.
\newblock \doi{10.1016/j.cpc.2009.09.018}

\bibitem[\protect\citeauthoryear{Sandberg, Drexler \& Ord}{2018}]{sandberg2018}
Sandberg A., Drexler E., Ord T., 2018, AcAau, 150, 24.
\newblock \doi{10.1016/j.actaastro.2018.04.030}

\bibitem[\protect\citeauthoryear{Troitskij}{1989}]{troitskij1989}
Troitskij V.~S., 1989, AcAau, 19, 875.
\newblock \doi{10.1016/0094-5765(89)90079-9}

\bibitem[\protect\citeauthoryear{Webb}{2002}]{webb2002}
Webb S., 2002, If the Universe Is Teeming with Aliens. Copernicus Books, New York

\end{thebibliography}

\end{document}
